% purpose: 读者应能回答"热量与水分在圆柱药材截面内分别沿什么方向传递、在表面受什么条件约束、两个场如何互相耦合"
% topology: geometry (轴对称剖面) + causal (双场耦合)
% reading_direction: coordinate-based（径向 r 由中心轴向外）
% node_roles: 物料域、对称轴、受热表面、热风边界层、温度场、含水率场
% edge_roles: 导热流、扩散质流、对流换热、表面对流传质、场间耦合
% visual_encoding: 实线=物料实体与热流；虚线=水分扩散；波浪线=对流边界层；点划线=对称轴；无颜色
% style_family: C1（纯黑白线稿，与全篇 mono 流程图一致）
% annotation_policy: minimal   result_policy: none   final_width: text   print_mode: grayscale
\documentclass[border=2pt]{standalone}
\usepackage{ctex}
\usepackage{amsmath}
\usepackage{tikz}
\usetikzlibrary{arrows.meta,calc,decorations.pathmorphing}
\begin{document}
\begin{tikzpicture}[
  font=\small,
  solidline/.style={line width=0.9pt},
  thinline/.style={line width=0.5pt},
  heatflow/.style={-{Stealth[length=4.4pt,width=3.2pt]},line width=0.9pt},
  massflow/.style={-{Stealth[length=4.4pt,width=3.2pt]},line width=0.7pt,densely dashed},
  couple/.style={-{Stealth[length=4pt,width=3pt]},line width=0.6pt},
  dimline/.style={{Stealth[length=3.6pt,width=2.8pt]}-{Stealth[length=3.6pt,width=2.8pt]},line width=0.5pt}
]
% ---- 半剖面：左侧为对称轴，右侧为受热表面。半剖面而非整圆，便于把径向坐标画成真正的一维轴。
\def\Rr{4.2}      % 表面半径位置
\def\Hh{3.1}      % 剖面半高

% 物料区域轮廓（实线=给定实体边界）
\draw[solidline] (0,-\Hh) -- (\Rr,-\Hh) -- (\Rr,\Hh) -- (0,\Hh);
% 对称轴（点划线=辅助线）
\draw[thinline,dash pattern=on 6pt off 2pt on 1pt off 2pt] (0,-\Hh-0.45) -- (0,\Hh+0.45);
% 抬 0.12：与右上两行标签的顶边取齐（那块两行、顶边比锚点高一整行）
\node[anchor=south,inner sep=1.5pt] at (0,\Hh+0.62) {对称轴};

% ---- 径向坐标轴（双箭头量度线，置于图形之外）
\draw[dimline] (0,-\Hh-0.95) -- (\Rr,-\Hh-0.95);
\node[anchor=north,inner sep=2pt] at (0,-\Hh-1.0) {$r=0$};
\node[anchor=north,inner sep=2pt] at (\Rr,-\Hh-1.0) {$r=R$};
\node[anchor=south,inner sep=2pt] at ({\Rr/2},-\Hh-0.92) {径向坐标 $r$};

% ---- 表面与边界层：波浪线表示热风侧对流边界层
\draw[thinline,decorate,decoration={snake,amplitude=1.1pt,segment length=6pt}]
      (\Rr+0.62,-\Hh) -- (\Rr+0.62,\Hh);
\node[anchor=south,align=center,inner sep=2pt] at (\Rr+0.62,\Hh+0.10) {热风侧\\对流边界层};

% ---- 热流：由热风穿过边界层进入表面，再向轴心导入（实线箭头）
% 三行热流等间距（1.50 个单位）：原来 1.20/1.80 的不等间距会让中段读成"漏了一项"。
% 内侧箭头尖端收在 2.80，与耦合竖箭头(x=2.45)留出净空，避免几条线在同一小区域挤作一团。
\foreach \y in {2.15,0.65,-0.85}{
  \draw[heatflow] (\Rr+1.62,\y) -- (\Rr+0.08,\y);
  \draw[heatflow] (\Rr-0.22,\y) -- (2.80,\y);
}
\node[anchor=west,align=left,inner sep=2pt] at (\Rr+1.72,2.15) {对流换热\\第三类边界条件};
\node[anchor=west,inner sep=2pt] at (\Rr+1.72,0.65) {表面吸热升温};
\node[anchor=west,inner sep=2pt] at (\Rr+1.72,-0.85) {向内导热};

% ---- 质流：水分自内向外扩散并在表面蒸发（虚线箭头）
\foreach \y in {-1.30,-2.35}{
  \draw[massflow] (2.80,\y) -- (\Rr-0.10,\y);
  \draw[massflow] (\Rr+0.10,\y) -- (\Rr+1.54,\y);
}
\node[anchor=west,inner sep=2pt] at (\Rr+1.72,-1.30) {水分向外扩散};
\node[anchor=west,align=left,inner sep=2pt] at (\Rr+1.72,-2.35) {表面对流传质\\第三类边界条件};

% ---- 域内两个场与它们的耦合（同一画面内，不另起文字栏）
\node[inner sep=2.2pt] at (1.55,1.62) {温度场 $T(r,t)$};
\node[inner sep=2.2pt] at (1.55,-1.62) {含水率场 $C(r,t)$};
% 两端各收 0.12 单位：Stealth 箭尖原先离两个场标签只有 1.6-2.3pt，看着像刺进文字。
% 标签 inner sep 加到 4pt：原 2pt 时首字距箭杆 2.3pt，属"文字贴线"。
\draw[couple] (1.80,1.26) -- node[left,inner sep=4pt,align=right]{升温增强\\水分扩散} (1.80,-1.26);
\draw[couple] (2.45,-1.26) -- node[right,inner sep=4pt,align=left]{失水改变\\热物性} (2.45,1.26);

% ---- 中心对称条件与内外差异（贴在轴侧，说明零通量与温升滞后）
\node[anchor=west,inner sep=2pt] at (0.10,-2.62) {轴心零通量};
\node[anchor=west,inner sep=2pt] at (0.10,2.72) {中心温升滞后于表面};
\end{tikzpicture}
\end{document}
